% Осень26!A15; нумерация внутри столбца без пустых строк.
% Источник: лекция 5, §14 (строка 487)
\ticket{13}{Лемма о приближенных значениях дифференцируемого отображения}
\label{ticket:13}

\materialnote{В лекции лемма названа «о почти локальной инъективности». Соответствие названию в списке требует сверки.}

\subsection*{Лемма о почти локальной инъективности}

\emph{(Лектор: «назову её нелепое название --- почти локальная инъективность»; в списке билетов под №13 фигурирует как «лемма о приближённых значениях дифференцируемого отображения» --- \textbf{Требует сверки.} соответствие названий; на лекции доказана оценка снизу, см. ниже.)}

\textbf{Мотивация.} По определению дифференцируемости \[F(x_0+h)-F(x_0)=F'(x_0)\,h+o\bigl(|h|\bigr),\] и мы хотим воспринимать разность слева как «выражение первого порядка по \(h\)» (в одномерном случае отклик линеен, коэффициент --- производная). Но если оператор \(F'(x_0)\) \textbf{вырожден}, он может занулять или сильно гасить отдельные приращения --- точки будут «склеиваться» или слишком сближаться, и оценка порядка \(|h|\) невозможна. Лемма: для \textbf{невырожденного} оператора всё хорошо --- «железно» существует константа, ограничивающая снизу.

\begin{lemma}

Пусть \(O\subset\mathbb R^m\) открыто, \(F:O\to\mathbb R^m\), \(x_0\in O\), \(F\) дифференцируема в точке \(x_0\) и \(\det F'(x_0)\ne0\). Тогда \[\exists\,c>0\ \exists\,\delta>0\ \forall h\in\mathbb R^m:\ |h|<\delta\ \Rightarrow\ |F(x_0+h)-F(x_0)|\ge c\,|h|.\] (Точка \(x_0\) «не склеивается с близкими».)

\end{lemma}

\subsubsection*{Доказательство}

\textbf{Шаг 1: частный случай --- линейное отображение .} Пусть \(F=A\) линейно. Тогда \(F'(x_0)=F=A\) при любом \(x_0\): действительно, из линейности \(F(x_0+h)-F(x_0)=F(h)=A h+0\), и годится представление дифференцируемости с «\(o\)-малым», тождественно равным нулю; по единственности производного оператора \(F'(x_0)=A\). (Пример: \(F(x)=10x\) из \(\mathbb R\) в \(\mathbb R\): \(F'=10\); для линейного отображения \(F'=F\) «в любой точке».)

Так как \(\det A\ne0\), существует обратный оператор \(A^{-1}\). Из оценки для обратного оператора (норма линейного оператора, 1-й семестр): \(\forall y\) \[|A^{-1}y|\le\|A^{-1}\|\cdot|y|,\] откуда при \(y=Ah\): \[|h|\le\|A^{-1}\|\cdot|A h|\qquad\Longrightarrow\qquad |A h|\ge\frac{1}{\|A^{-1}\|}|h|.\] Но для линейного \(F\): \[|F(x_0+h)-F(x_0)|=|Ah|,\] и константа \(c:=1/\|A^{-1}\|\) подходит при любом \(\delta\) («линейному оператору без разницы все эти фокусы с областями определения»).

\textbf{Шаг 2: общий случай .} Выпишем определение дифференцируемости в \(x_0\), выделив бесконечно малую: \[F(x_0+h)-F(x_0)=F'(x_0)h+\alpha(h)\,|h|,\qquad \alpha(h)\to0\ (h\to0).\] По неравенству треугольника и шагу 1 (с \(A=F'(x_0)\), \(c_0:=1/\|(F'(x_0))^{-1}\|\)): \[|F(x_0+h)-F(x_0)|\ \ge\ |F'(x_0)h|-|\alpha(h)|\,|h|\ \ge\ \bigl(c_0-|\alpha(h)|\bigr)|h|.\] Выберем \(\delta>0\) так, чтобы \(|\alpha(h)|\le c_0/2\) при \(|h|<\delta\). Тогда \[|F(x_0+h)-F(x_0)|\ge\frac{c_0}{2}\,|h|,\qquad c:=\frac{c_0}{2}.\qquad\blacksquare\]

\subsubsection*{Замечание: глобальной инъективности нет}

Из \(\det F'(x)\ne0\) \textbf{во всех} точках не следует инъективность \(F\) на всём \(O\): «точка не склеивается с близкими, но может склеиться с далёкими».

\textbf{Контрпример.} В координатах \((u,v)\) рассмотрим \[F(u,v)=\bigl(u^2-v^2,\ 2uv\bigr)\] --- возведение комплексного числа \(u+iv\) в квадрат. Матрица Якоби \[F'(u,v)=\begin{pmatrix}2u & -2v\\ 2v & 2u\end{pmatrix},\qquad \det F'(u,v)=4(u^2+v^2),\] т.е. \(\det F'\ne0\) всюду на \(\mathbb R^2\setminus\{(0,0)\}\). Но \(F(1,0)=F(-1,0)=(1,0)\) --- инъективности нет.
